\section{Multiprocessing and networking in SMC}
Any robotics software has to be capable of parallel computation.
Simply put, you don't want your controller to stall because the processor
is busy running a visualization, or waiting on a camera update.
Normally, this would be handled with multithreading,
but since Python still has the Global Interpreter Lock (GIL),
all Python threads run on a single processor core.

For this reason, we are forced to implement parallelism via processes.
Fortunately, since SMC is a library for quick-and-dirty implementation
for testing and throwing away robotics experiments, we can get away with
spawning all required processes in the beginning, and rest assured
that they will run in the 5 minutes our experiments last
(if you want serious code for serious applications, why are you looking at a Python library?).

Interprocess communication (IPC) in SMC is designed to be as lean and as quick as possible.
In other words, \textbf{we always aim to minimize the amount of data transferred
between processes, and to always have the last piece of data available},
irrespective of the actual protocol used.

There are 3 IPC techniques in SMC: 
\begin{itemize}
\item \textbf{queues} - for asynchronous data transfer for processes
		running at substantially different frequencies,
		ex. for visualization
\item \textbf{shared memory} - for instantaneous data access
		for processes running at comparable frequencies in the same machine
\item \textbf{sockets} - for processes running at different machines,
		with serialization from protobuf (allows for heterogeneous software stacks)
\end{itemize}

From the user point of view, IPC techniques are hidden
behind basic multi-process interaction paradigms, which
are implemented in a dedicated \fbox{\texttt{ProcessManager}}.
The main process (the one spawned when running an example script) is reserved
for control. Other processes are spawned by passing a function
to be run in a separate process
to a particular \fbox{\texttt{ProcessManager}}, which will be introduced next.

\subsection{ProcessManager}
\fbox{\texttt{ProcessManager}} is a wrapper class around the function
you want to run in a separate process.
That function is referred to as \fbox{\texttt{side\_function}}.
It's sole purpose is to spawn the process and the selected IPC method
for you. It's \textbf{your job} to declare which data you're exchanging
to actually send and receive said data, in the same manner as for logging
and plotting.

The following multi-process interaction paradigms are supported in SMC:
\begin{itemize}
\item \fbox{\texttt{ConsumerProcess}} - the side function only takes
	in data from the control process via queue (ex. visualizer, plotter)
\item \fbox{\texttt{ProducerProcess}} - the side function only provides 
	data to the control process via queue (ex. camera outputs)
\item \fbox{\texttt{ProducerConsumerProcess}} - combines
	functions of the consumer and producer processes
\item \fbox{\texttt{CollaborativeProcess}} - like producer-consumer
	process, but the communication is implemented via shared memory.
	Allows for nearly-instantaneous communication,
	ex. required in path planning (provide current position,
	obtain current path)
\end{itemize}

The producer processes require an \fbox{\texttt{init\_value}} argument,
which will be return before the first producer output.
The consumer processes require an \fbox{\texttt{init\_command}}
argument to begin computation at process spawn time, instead
of when the controller starts.
Additional arguments to the side function are to be provided 
via \fbox{\texttt{partial}}-ing.

\paragraph{Using ProcessManagers}
The process with your function will be started 
immediately upon construction of a ProcessManager object. 
You then interact with the process with the following functions called
from the ProcessManager object:
\begin{itemize}
\item \fbox{\texttt{sendCommand(command)}} - sends the command/message
		to a consumer process 
\item \fbox{\texttt{getData()}} - obtain the latest information
	from the producer process - \textbf{always}
		returns the latest received data
\end{itemize}
Depending on the underlying IPC used, the side functions will need
to receive different IPC objects, which then determine
how \fbox{\texttt{sendCommand }} and \fbox{\texttt{getData}} work.
The simplest way to navigate this is to look at existing examples
and replicating the procedure therein.

\subsection{Multiprocessing examples}
\paragraph{Producer-consumer examples}
Producer-consumer is multiprocessing lingo for a communication
relationship where one process ``produces'' information (ex. messages, tasks,
commands) and the other process ``consumes'' that information.
In other words, the communication goes one way only.
Another popular name for this communication relationship is publisher-subscriber.

A hands-on example can be found in
\fbox{\texttt{examples/vision/camera\_no\_lag.py}}
In this case control is the consumer, and the side process is the producer.
The side process obtains images from a camera, processes them,
and puts the result in a queue. The results are then made available
in the control loop. This example serves only to provide
the structure of this would work, but there is no actual task.
An example of how this setup is used 
for visual servoing can be found in
\fbox{\texttt{examples/planar\_pushing\_via\_visual\_servoing}.}

This communication relationship can also be found throughout the codebase.
The control loop is a producer for both the plotter
and the visualizer. Both use queues. 
In path following control from a planner, the path planner consumes
the current position of the robot, and produces the path to be followed.
Since this needs to be as close to instantaneous as possible,
shared memory is used for both pieces of information.


\subsection{Networking}
Working with networking requires a basic operational knowledge of computer
networks (what an IP address is etc). This is unavoidable
because it is very likely you'll have to set up and debug the network yourself.
You are referred to 
\href{https://www.youtube.com/watch?v=0Hhaf4q3eqI}{an introductory video} 
if you have no experience with networking.

More info on networking in SMC can be found in 
\newline
\fbox{\texttt{smc/multiprocessing/networking/notes.md}}.
Examples of networking in SMC can be found in 
\newline
\fbox{\texttt{examples/networked\_clik\_client}}
One example covers connecting 2 SMC clients, and the other connecting 
an SMC client with a ROS1 client.




\section{Multiprocessing in ROS2}
If you are a ROS2 user, just use ROS2 for multiprocessing.
The bare-bones SMC implementation serves primarily to avoid 
ROS2 as a dependency (as this would be the only thing in it that's used).
RCL automatically assigns different IPC protocols to
its publisher-subscriber structures, and you do not have to think
about what is ``under the hood.''
If you have to use ROS2 for interfacing with your robot,
you will have to learn how it works to access all features of your robot.
In this case you are referred to \href{https://docs.ros.org/en/rolling/index.html}{ROS2 documentation},
and are advised to use only control from SMC.
This requires a specific pre-set argument list, and a particular node structure
to work. More on this in Appendix \ref{sec-ros2}.
